\documentclass{article} % For LaTeX2e
\usepackage{iclr2027_conference,times}

% Optional math commands from https://github.com/goodfeli/dlbook_notation.
% \input{math_commands.tex}   % not used in this paper

\usepackage{hyperref}
\usepackage{url}
\usepackage{booktabs}
\usepackage{graphicx}
\usepackage{amsmath}
\usepackage{multirow}

\title{Knowing the Answer Is Not Enough: A Vietnamese Medical\\
Multiple-Answer Benchmark with Derived Ground Truth}

% Authors must not appear in the submitted version. They should be hidden
% as long as the \iclrfinalcopy macro remains commented out below.
% Non-anonymous submissions will be rejected without review.
\author{Anonymous authors\\
Paper under double-blind review}

\newcommand{\fix}{\marginpar{FIX}}
\newcommand{\new}{\marginpar{NEW}}

\iclrfinalcopy % Uncomment for camera-ready version, NOT for submission.

\begin{document}

\maketitle

\begin{abstract}
Clinical questions often admit several correct answers, and knowing how many is
itself part of competence; yet medical question answering is evaluated almost
exclusively in the single-answer multiple-choice format, which hands the model that
count for free. To our knowledge no Vietnamese medical benchmark tests the
multiple-answer setting, and the two obvious ways to build one both fail: prompting a
model to author items makes the benchmark a measure of agreement with the generator,
while expert authoring does not scale to tens of thousands of items. We instead
\emph{derive} multiple-answer items from single-answer exam questions whose keys are
already known, using four rule-based transformations over recurring structures in
Vietnamese exam writing---negated stems, ``all of the above'' keys, letter-referencing
keys, and options that are themselves comma-separated lists---so that every
correct-answer set follows deterministically from the source key; a language model only
screens candidates, rewrites stems and writes distractors, never deciding which options
are correct, and an independent audit round re-examines every assembled item in both
directions. From 17 exam files (160,194 items; 122,017 after global deduplication) we
release 9,368 multi-answer items and an 11,368-item mixed set that conceals the answer
count, each paired with an index-aligned single-answer control that separates format
effects from difficulty. Across nine models exact match spans 1.8--61.5\% and 20.3\% of
items are solved by none. The binding constraint is answer-set calibration rather than
knowledge: GPT-4o-mini's selection contains the correct option for 90.5\% of the
single-answer items, the highest of all models evaluated, yet it answers only 13.2\%
exactly, rising to 63.5\% once told the count. Enabling reasoning on a fixed model does
not repair this deficit and slightly worsens it, whereas successive model generations
improve it sharply---evidence that answer-set calibration is a learned capability rather
than a product of inference-time deliberation.
\end{abstract}

\section{Introduction}

Medical question answering benchmarks are near-universally single-answer multiple
choice. MedQA \citep{jin2021medqa}, MedMCQA \citep{pal2022medmcqa}, PubMedQA
\citep{jin2019pubmedqa} and the medical subsets of MMLU \citep{hendrycks2021mmlu} all
present a stem, a handful of options, and exactly one key. The format is convenient:
scoring is unambiguous and a random-guess baseline is trivial to compute. But it embeds
an assumption that does not hold in clinical practice. Real diagnostic and management
questions frequently admit several defensible answers, and part of competence is knowing
how many.

Standard multiple-choice evaluation gives that information away. When the prompt states
that exactly one option is correct, a model that has narrowed the field to two candidates
and guesses between them is scored identically to a model that is certain. Worse, a model
that has correctly identified every true statement is \emph{penalised} for reporting more
than one. The format measures option ranking, not answer-set selection.

We argue that answer-set selection is separately measurable, separately deficient in
current models, and cheap to evaluate provided the benchmark is constructed carefully.

\paragraph{The construction problem.} The obvious way to build a multiple-answer
benchmark is to prompt a strong language model to write the items. This is circular: the
resulting benchmark measures agreement with the generating model, and any systematic
error in the generator silently becomes ground truth. The alternative---expert
authoring---does not scale to the tens of thousands of items needed for stable
per-subgroup estimates.

Our approach is neither. We observe that large Vietnamese medical exam banks already
contain multiple-answer items in disguise. A question of the form \emph{``All of the
following are risk factors for X, EXCEPT: \ldots''} with key D asserts that options A, B
and C are all true. A question whose key is \emph{``All of the above''} asserts the same
of every other option. A question whose four options are each comma-separated lists of
findings draws from a shared pool of concepts, and the key identifies which pool members
belong. In each case the multi-answer ground truth is \emph{derived} from the existing
answer key by a deterministic rule rather than authored, and the source item's
correctness guarantee transfers.

This shifts the role of the language model from author to filter. In our pipeline the
model screens candidates for whether a transformation is semantically sound, rewrites
negated stems into affirmative form, and writes distractor text where a transformation
removes options. It never selects the correct answers. The distinction matters: a
generated distractor that turns out to be true is a \emph{detectable} defect---an
independent audit round flags it---whereas a generated \emph{key} that is wrong is
indistinguishable from ground truth.

\paragraph{The mixed-format design.} To measure cardinality calibration we need items
where the model does not know the answer count a priori. We therefore release a second
benchmark in which 2,000 single-answer items---drawn from the same source bank and
disjoint from the multi-answer items by three independent criteria---are interleaved
with the multi-answer pool, and the prompt states only that a question ``may have one or
more correct answers.'' Because those 2,000 items also appear verbatim in the paired
single-answer control, the same model can be scored on the same questions with and
without the cardinality hint. The difference is a direct measurement of calibration cost.

\paragraph{Contributions.}
\begin{enumerate}
\item A construction method that yields multiple-answer items at scale with
      \emph{derived} rather than generated ground truth, together with a two-round LLM
      protocol in which the model is structurally prevented from choosing answers
      (Section~\ref{sec:construction}).
\item Two paired Vietnamese medical benchmarks---9,368 multi-answer items and an
      11,368-item mixed set---each shipped with an index-aligned single-answer control
      that isolates format effects from difficulty (Section~\ref{sec:sets}).
\item An evaluation of nine models showing that answer-set cardinality calibration, not
      knowledge, is the binding constraint, and that this deficit is not repaired by
      inference-time reasoning but improves sharply across model generations
      (Section~\ref{sec:results}).
\end{enumerate}

\section{Related Work}

\paragraph{Medical question answering benchmarks.} MedQA \citep{jin2021medqa} draws from
professional medical board examinations; MedMCQA \citep{pal2022medmcqa} covers Indian
medical entrance exams across 21 subjects; PubMedQA \citep{jin2019pubmedqa} poses
yes/no/maybe questions grounded in research abstracts. Med-PaLM
\citep{singhal2023clinical} established the aggregate-benchmark framing that most
subsequent work follows. All of these evaluate in a format where the number of correct
answers is fixed and known to the model in advance. Our benchmark removes that
assumption.

\paragraph{Limitations of the single-answer format.} The format itself is a known source
of measurement error: models exhibit systematic position and token biases as multiple
choice selectors rather than reading options on their merits \citep{zheng2024selectors},
and MMLU-Pro \citep{wang2024mmlupro} responds by enlarging the option set and raising
reasoning difficulty. Both of these keep the number of correct answers at one. We vary
that number instead.

\paragraph{Benchmark construction with language models.} Using a strong model to generate
evaluation items is now common practice \citep{wang2023selfinstruct}, and its circularity
is a recognised hazard: the resulting scores partly measure agreement with the generator,
and models demonstrably favour their own generations when placed in an evaluative role
\citep{zheng2023judging,panickssery2024selfpreference}. Our derivation approach sidesteps
this by constraining the model to roles where its errors are detectable by an independent
pass rather than baked into the key.

\paragraph{Non-English medical evaluation.} Medical LLM evaluation remains concentrated
in English and Chinese, the latter served by large exam-derived benchmarks such as CMB
\citep{wang2024cmb} and CMExam \citep{liu2023cmexam}. Vietnamese medical question
answering is comparatively under-served: existing resources target dialogue intent
\citep{huy2021vimq} and abstractive answering \citep{tran2024vimedaqa}, while the broad
Vietnamese LLM evaluation suites \citep{truong2024vietnamese} are not medical-specific.
We are not aware of a prior Vietnamese benchmark in the multiple-answer setting.

\paragraph{Multi-label and set-valued prediction.} Predicting a set rather than a single
label is a long-studied problem \citep{tsoumakas2007multilabel}, and exact set match is
the natural strict metric; conformal methods likewise target prediction sets whose size
adapts to difficulty \citep{romano2020classification}. What is new here is not the metric
but the observation that a model's failure to control the \emph{size} of its predicted
set dissociates cleanly from its knowledge of the set's contents---a failure of
calibration in the sense of \citet{guo2017calibration}, but over set cardinality rather
than over confidence---and that this dissociation is invisible under single-answer
evaluation.

\section{Benchmark Construction}
\label{sec:construction}

\subsection{Source corpus and deduplication}

The source is 17 exported files from a Vietnamese medical exam bank, comprising 160,194
items in the schema \texttt{\{question, answer\_optionA..E, correctAnswer, medicalTopic,
difficultLevel\}}.

Deduplication must be global rather than per-file, because the exports overlap heavily.
We remove exact duplicates by SHA-1 over the normalised stem and option set, and
order-invariant duplicates by comparing the stem against the sorted option multiset.
Together these remove \textbf{30,913 items (19.3\%)}---24,519 exact and 6,394 reordered.
A further 7,264 items are dropped for having fewer than four options or an unparseable
key, leaving \textbf{122,017}.

One artefact deserves note because it silently leaks the answer. In 5,889 items every
option text is prefixed with its own letter (\texttt{"A. A. Chronic obstructive pulmonary
disease"}). Left in place, options carried over from the source would bear the prefix
while model-written distractors would not, letting a model identify the correct options
without reading them. We strip the prefix when at least two options carry a
position-matching marker, and do so \emph{before} computing the deduplication hash---which
also merges export variants differing only by this prefix, raising detected duplicates
from 29,797 to 30,913.

\subsection{Four transformation families}

\begin{table}[t]
\caption{The four transformations. Each converts a single-answer item into a
multiple-answer item whose key follows deterministically from the source key.}
\label{tab:families}
\begin{center}
\begin{small}
\begin{tabular}{llr}
\toprule
\textbf{Family} & \textbf{Derivation} & \textbf{Candidates} \\
\midrule
\textsc{negation} & stem asks for the single false option; the other three & 4,729 \\
                  & become the correct set, and the stem is rewritten affirmatively & \\
\textsc{pool}     & every option is a comma-separated list; list items merge into & 3,048 \\
                  & a shared pool, and the key's items become the correct set & \\
\textsc{meta-all} & key is ``All of the above''; the meta option is removed, the & 1,102 \\
                  & rest are all correct, and the model writes two distractors & \\
\textsc{meta-letters} & key is ``Both A and B''; the correct set is read off the & 900 \\
                  & letters, remaining options are human-authored distractors & \\
\midrule
& \textbf{Total} & \textbf{9,779} \\
\bottomrule
\end{tabular}
\end{small}
\end{center}
\end{table}

Table~\ref{tab:families} summarises the four transformations. Mechanical filters run
before any model call: decimal commas are protected before list splitting; option pools
that are predominantly numeric are rejected; pools whose options do not share a concept
vocabulary are rejected by an overlap threshold; ordered-sequence and ranking items are
detected by repeated items within an option and by identical item sets across options;
blank-slot stems and junk options are removed; and near-duplicate stems are removed by
blocked \texttt{SequenceMatcher} comparison.

For the \textsc{pool} family we accept pools of up to eight items and then randomly trim
to five options, retaining every correct item and dropping only distractors. Hard-capping
at five during filtering would have retained 1,280 candidates; trimming afterwards retains
3,048 with the same final shape.

\subsection{Two-round LLM protocol}

All model calls use \texttt{Qwen3-32B-FP8} under vLLM with structured outputs and thinking
disabled.

\paragraph{Round 1 (9,779 calls)} is family-specific: verify that a negation item is
genuinely negated and rewrite its stem affirmatively; write two distractors for the two
meta families; and for \textsc{pool}, judge whether the items are independently meaningful
and flag any \emph{distractor that is in fact also correct}. Only removal is permitted---the
model may delete a pool item, never add one to the key. Of 9,779 responses, 9,778 parsed;
9,368 items survived construction.

\paragraph{Round 2 (9,368 calls)} is uniform across families: the assembled item is
presented with its current correct/incorrect labels, and the model is asked, in both
directions, which labelled-correct options are actually wrong and which labelled-incorrect
options are actually right. This is the only stage that can catch \emph{semantic}
defects---no syntactic rule knows that a particular distractor is a true statement.

Two failure modes in the audit output required handling. First, the model conflates the
two index lists: 26.5\% of indices in the ``labelled-correct but actually wrong'' array
point at options already labelled incorrect. We discard such self-contradictory verdicts
rather than treat them as evidence; counting them would reject 1,229 items unnecessarily.
Second, the audit is not position-neutral---it accepts items whose single incorrect option
sits last more often, raising the share of incorrect options at position D from 21.6\% to
24.0\% ($\approx 5\sigma$). We therefore re-shuffle option positions \emph{after} the audit
gate, which is safe because surviving items have no flagged label errors.

\subsection{The two released benchmarks}
\label{sec:sets}

\textbf{Set A} is the 9,368 multi-answer items, whose correct-answer sets contain two
options (1,901 items), three (7,102), or four (365). \textbf{Set B} adds 2,000
single-answer items drawn verbatim from the source bank, disjoint from Set A by item hash,
by \texttt{(source file, index)} pair, and by near-duplicate stem matching against both the
multi-answer items and their originals; they are spread evenly at 117--118 items per source
file. Items whose key is a meta option are excluded, since those are multi-answer items in
disguise. Set B therefore contains 11,368 items: 4,573 \textsc{negation}, 2,810
\textsc{pool}, 1,092 \textsc{meta-all}, 893 \textsc{meta-letters}, and 2,000 single-answer.

Each set ships an \textbf{index-aligned single-answer control}: record $i$ of the control is
the verbatim source question for record $i$ of the benchmark. Differences between the two
are therefore attributable to format rather than to item difficulty.

Option positions are shuffled twice with seeds derived from the item hash, and item order
once with a fixed seed. In Set B the incorrect-option position distribution is
A 4,697 / B 4,712 / C 4,918 / D 4,908 (and E 1,858, which occurs only in five-option
items), showing no positional leak.

\paragraph{A caveat we report rather than hide.} The single-answer controls inherit a
strong position bias from the source: 41.5\% of Set A control keys are option D, because
the \textsc{meta-all} and \textsc{meta-letters} families select items whose key is a meta
option, and meta options are almost always written last (89.7\% and 84.3\% D
respectively). An always-D baseline scores 41.5\% on that control. Aggregate control
scores are therefore inflated; the paired per-item comparison, which is what we use, is
unaffected.

\section{Experimental Setup}
\label{sec:setup}

We evaluate nine models. Four are open-weight (Qwen3-0.6B and Qwen3-4B,
\citealp{yang2025qwen3}; Gemma3-1B and Gemma3-4B, \citealp{gemmateam2025gemma3}) and run
under vLLM \citep{kwon2023vllm} on a single A10 24GB GPU; five are API models
(GPT-4o-mini, \citealp{openai2024gpt4o}; DeepSeek-v4-Flash, in the model family of
\citealp{deepseekai2024v3}; Gemini-2.5-Flash, \citealp{comanici2025gemini25};
Claude-Sonnet-5; Gemini-3.6-Flash) accessed through a common OpenAI-compatible interface.
Exact endpoint identifiers and provider versions for every API model are recorded in the
released harness configuration.

Decoding uses greedy sampling where the provider permits it, and output is constrained to a
JSON schema whose enum is exactly the option letters of that item. Constraining the output
is essential for comparability: the open models run under grammar-constrained decoding, so
an unconstrained API model would effectively be scored on a different task. The primary
metric is \textbf{exact set match}; Jaccard and micro-F1 are reported as partial-credit
secondaries.

\paragraph{Reasoning conditions.} Eight of the nine models are evaluated with reasoning
disabled, verified by reading back the provider's reported reasoning-token count and
confirming it is zero. Gemini-3.6-Flash cannot be evaluated this way---the endpoint returns
\texttt{HTTP 400: "Reasoning is mandatory for this endpoint and cannot be disabled"}---so its
scores are reported separately and never pooled with the matched-condition group.
Section~\ref{sec:ablation} reports a controlled reasoning ablation on a single model.

\paragraph{Models we could not evaluate under matched conditions.} Several frontier
endpoints (\texttt{gpt-5-mini}, \texttt{gemini-2.5-pro}, \texttt{gemini-3.5-flash},
\texttt{grok-4.6}) also mandate reasoning. One further model, \texttt{grok-4.3}, runs
correctly but was excluded on cost: it tokenises accented Vietnamese at 1.09 characters per
token against Gemini's 3.11, making a full sweep roughly three times more expensive than its
per-token price implies. \textbf{Tokeniser efficiency on the target language is a
first-order cost factor for non-English benchmarks and is not visible in published
pricing}; measured input cost per item varied by $3.6\times$ across providers for identical
prompts.

\section{Results}
\label{sec:results}

\subsection{Overall accuracy}

\begin{table}[t]
\caption{Exact set match (\%). Set A is multi-answer only; Set B interleaves 2,000
single-answer items. Controls contain the verbatim source questions at matching indices.
Gemini-3.6-Flash is listed separately because reasoning cannot be disabled on that
endpoint; it was evaluated on Set B only, as was Claude-Sonnet-5.}
\label{tab:main}
\begin{center}
\begin{small}
\begin{tabular}{lrrrr}
\toprule
\textbf{Model} & \textbf{Set A} & \textbf{Set B} & \textbf{Control A} & \textbf{Control B} \\
\midrule
\multicolumn{5}{l}{\emph{Reasoning disabled (matched conditions)}} \\
Gemma3-1B          &  2.1 &  1.8 & 37.5 & 37.0 \\
Qwen3-0.6B         &  3.6 &  3.3 & 40.6 & 39.5 \\
Gemma3-4B          & 20.5 & 16.3 & 50.2 & 49.3 \\
Qwen3-4B           & 24.7 & 21.2 & 55.7 & 55.5 \\
GPT-4o-mini        & 37.5 & 33.9 & 67.0 & 66.2 \\
DeepSeek-v4-Flash  & 45.1 & 43.3 & 72.8 & 72.2 \\
Gemini-2.5-Flash   & 51.9 & 50.0 & 77.6 & 77.1 \\
Claude-Sonnet-5    &  --- & \textbf{58.0} & --- & --- \\
\midrule
\multicolumn{5}{l}{\emph{Reasoning mandatory (not comparable to the group above)}} \\
Gemini-3.6-Flash   &  --- & \textbf{61.5} & --- & --- \\
\bottomrule
\end{tabular}
\end{small}
\end{center}
\end{table}

Table~\ref{tab:main} reports the main results. Among models scored under matched
conditions, the ranking is identical across all four datasets, so conclusions do not depend
on which set is reported. The spread on Set B is $32\times$ (58.0 against 1.8).
Claude-Sonnet-5 leads the matched group by 8.0 points over Gemini-2.5-Flash; a paired
McNemar test gives 1,746 items solved only by Sonnet against 831 solved only by Gemini
($p \approx 10^{-72}$).

The two sub-2B models score below their own single-answer controls by more than an order of
magnitude. They do not merely perform poorly in the multiple-answer format---they fail at
it.

\subsection{Cardinality calibration dissociates from knowledge}
\label{sec:cardinality}

\begin{table}[t]
\caption{The 2,000 single-answer items inside Set B, scored twice: once where the prompt
does not reveal the answer count, and once in the control where it states there is exactly
one. ``Key $\subseteq$ selection'' is the fraction of items whose correct option appears
anywhere in the model's selected set. Dashes mark runs not performed.}
\label{tab:cardinality}
\begin{center}
\begin{small}
\begin{tabular}{lrrrrr}
\toprule
\textbf{Model} & \textbf{Told count} & \textbf{Not told} & \textbf{$\Delta$} &
\textbf{Selects exactly 1} & \textbf{Key $\subseteq$ selection} \\
\midrule
Qwen3-0.6B        & 33.1 &  1.2 & $-31.9$ &  1.9 & 56.9 \\
Gemma3-1B         & 34.8 &  0.1 & $-34.8$ &  0.1 & 54.0 \\
Gemma3-4B         & 45.0 &  4.7 & $-40.4$ &  7.4 & 71.6 \\
Qwen3-4B          & 54.1 & 12.5 & $-41.6$ & 16.6 & 71.7 \\
GPT-4o-mini       & 63.5 & 13.2 & $\mathbf{-50.2}$ & 15.0 & \textbf{90.5} \\
DeepSeek-v4-Flash & 67.4 & 41.4 & $-26.0$ & 52.1 & 82.0 \\
Gemini-2.5-Flash  & 74.4 & 50.9 & $-23.5$ & 63.2 & 83.2 \\
Claude-Sonnet-5   &  --- & 63.6 &  ---    & 77.8 & 83.9 \\
Gemini-3.6-Flash  &  --- & 73.2 &  ---    & \textbf{85.1} & 87.0 \\
\bottomrule
\end{tabular}
\end{small}
\end{center}
\end{table}

Table~\ref{tab:cardinality} isolates the central finding. GPT-4o-mini is the clearest case.
Its selected set contains the correct answer for 90.5\% of items---the highest of every
model we evaluated, above both Gemini-2.5-Flash and Claude-Sonnet-5---yet it answers only
13.2\% exactly, because it selects a single option in just 15.0\% of cases. It holds more
of the required knowledge than models that score four times higher, and converts far less
of it into answers. The capability gap is not knowledge; it is knowing when to stop.

Averaged over all of Set B, where the mean key size is 2.51, mean selection size ranges
from 2.01 (Gemma3-1B) to 3.06 (GPT-4o-mini). Gemma3-1B selects exactly two options for
11,315 of 11,368 items (99.5\%) regardless of the item---it is effectively blind to
cardinality. Three distinct failure modes are therefore visible: blindness to cardinality
in the smallest models, systematic over-selection in GPT-4o-mini, and residual
mis-calibration in the strongest models.

\subsection{Source difficulty labels are valid}

\begin{table}[t]
\caption{Exact set match (\%) on Set B by source difficulty label. The label is written by
the exam authors in 15 surface forms, which we normalise to four levels, and is never
consulted during construction.}
\label{tab:difficulty}
\begin{center}
\begin{small}
\begin{tabular}{lrrrrrrrrr}
\toprule
\textbf{Level} & \textbf{$n$} & \textbf{Gem3-1B} & \textbf{Qw-0.6B} & \textbf{Gem3-4B} &
\textbf{Qw-4B} & \textbf{4o-mini} & \textbf{DeepSeek} & \textbf{Gem-2.5} & \textbf{Son-5} \\
\midrule
1 Easy      & 2,405 & 2.0 & 4.2 & 17.5 & 26.2 & 35.6 & 49.8 & 57.0 & \textbf{65.7} \\
2 Moderate  & 7,496 & 1.7 & 3.1 & 16.5 & 21.0 & 34.6 & 42.9 & 49.6 & \textbf{57.1} \\
3 Hard      & 1,428 & 1.9 & 2.9 & 13.0 & 14.4 & 27.7 & 34.7 & 41.0 & \textbf{50.4} \\
4 Very hard &    39 & 2.6 & 2.6 &  7.7 & 12.8 & 12.8 & 25.6 & 23.1 & \textbf{33.3} \\
\bottomrule
\end{tabular}
\end{small}
\end{center}
\end{table}

Table~\ref{tab:difficulty} shows that every sufficiently capable model declines
monotonically across levels 1--3. The two sub-2B models do not, which is consistent with
their failure being one of formatting rather than comprehension. Level 4 contains 39 items
and should not be read as a trend.

That an independently-authored difficulty annotation predicts accuracy for every capable
model is evidence that the derived items retain the semantic content of their sources. This
is a validity signal that does not depend on any assumption about our construction rules.

\subsection{How much headroom remains}
\label{sec:headroom}

Accumulating models from weakest to strongest, the fraction of Set B solved by \emph{no}
model falls as follows:

\begin{center}
\begin{small}
Gemma3-1B 98.2 $\to$ +Qwen3-0.6B 96.4 $\to$ +Gemma3-4B 81.7 $\to$ +Qwen3-4B 69.8 $\to$
+GPT-4o-mini 54.7 \\
$\to$ +DeepSeek 41.5 $\to$ +Gemini-2.5-Flash 33.3 $\to$ +Claude-Sonnet-5 25.8 $\to$
+Gemini-3.6-Flash \textbf{21.1}
\end{small}
\end{center}

Counting the reasoning-enabled run of Section~\ref{sec:ablation} as an additional system
lowers the residual to \textbf{20.3\%}. The four open models alone leave 69.8\% unsolved;
adding five API models removes two thirds of that.

The residual 20.3\% is an upper bound on the union of genuinely hard items and items whose
inherited key is wrong. The audit round rejected 1,570 constructed items for having a
labelled-correct option that is actually false, so wrong inherited keys are known to be
present at a non-trivial rate. Separating the two components requires manual review of a
sample, which we have not done.

\subsection{Reasoning does not repair calibration; newer generations do}
\label{sec:ablation}

To ask whether inference-time deliberation closes the gap, we re-ran the leading
matched-condition model at the time of the sweep---Gemini-2.5-Flash---with reasoning
enabled at \texttt{effort: low}, holding weights, prompt and output constraints fixed so
that the reasoning budget is the only difference.

\begin{table}[t]
\caption{Controlled reasoning ablation on a single model (Gemini-2.5-Flash, full Set B),
and the same calibration measure across three model generations evaluated on the 2,000
single-answer items.}
\label{tab:ablation}
\begin{center}
\begin{small}
\begin{tabular}{lrr}
\toprule
 & \textbf{Reasoning off} & \textbf{Reasoning low} \\
\midrule
Set B, exact match            & 50.0 & \textbf{51.0} \\
Mean options selected (key: 2.51) & \textbf{2.48} & 2.54 \\
Matches key size              & \textbf{61.1} & 60.0 \\
2,000 single-answer items     & \textbf{50.9} & 50.0 \\
Selects exactly one option    & \textbf{63.2} & 58.8 \\
Reasoning tokens per item     & 0 & 480 \\
\bottomrule
\end{tabular}

\vspace{1.5em}

\begin{tabular}{lrrr}
\toprule
\textbf{Selects exactly one option} & Gemini-2.5-Flash & Claude-Sonnet-5 & Gemini-3.6-Flash \\
\midrule
(2,000 single-answer items) & 63.2 & 77.8 & \textbf{85.1} \\
\bottomrule
\end{tabular}
\end{small}
\end{center}
\end{table}

The upper half of Table~\ref{tab:ablation} shows the result. Exact match rises by one
point, from 50.0\% to 51.0\%. A paired McNemar test finds 1,060 items correct only with
reasoning against 945 correct only without ($\chi^2 = 6.5$, $p = 0.011$): the difference is
statistically real, but the net gain of one point is the residue of 2,005 items changing
answer in both directions.

More informative is that \emph{every} calibration measure moves the wrong way. Key-size
match falls from 61.1\% to 60.0\%, accuracy on the 2,000 single-answer items falls from
50.9\% to 50.0\%, and the rate of selecting exactly one option falls from 63.2\% to 58.8\%.
Spending 480 reasoning tokens per item buys a marginal gain in knowledge and a measurable
loss in knowing when to stop. We verified that this is not an artefact of too small a
budget: on a 100-item probe, \texttt{effort: low} and \texttt{effort: medium} both scored
39\%.

The lower half of Table~\ref{tab:ablation} contrasts this with what changes across model
generations. On the same measure that reasoning degrades, successive generations improve
sharply: 63.2\% $\to$ 77.8\% $\to$ 85.1\%. The two halves together support a specific
claim: \textbf{answer-set calibration is a learned capability, not a product of
inference-time deliberation}. It is exactly the kind of capability that single-answer
evaluation cannot see, because single-answer evaluation supplies the answer for free.

\section{Limitations}

\paragraph{Distractors are model-written and not expert-verified.} 1,985 items in Set B
contain distractors authored by Qwen3-32B. They are filtered for length, meta-phrasing and
near-duplication against the correct options, and pass the round-2 audit, but no clinician
has reviewed them. Items carry a \texttt{qualityFlag} allowing them to be excluded.

\paragraph{Ground truth is inherited, and the source is imperfect.} Our derivations
preserve the source answer key, including its errors. A concrete example found during
inspection: an item on short-acting insulin lists ``duration 24 hours'' as correct, where
the correct value is 6--8 hours. The audit round removes some such items but cannot remove
all.

\paragraph{The single-answer controls have a position bias} (41.5\% key-D on Control A),
inherited from the meta-answer families. Aggregate control accuracy is inflated; the paired
per-item analysis is not affected.

\paragraph{Not every model was scored on every set.} Claude-Sonnet-5 and Gemini-3.6-Flash
were evaluated on Set B only, so their rows in Tables~\ref{tab:main} and
\ref{tab:cardinality} are partial, and no told-count contrast is available for them.

\paragraph{Greedy decoding was not available on every endpoint.} Some frontier models no
longer accept a temperature parameter, so their runs use provider-default sampling. Scores
for those models carry run-to-run variance that the grammar-constrained open-weight runs do
not.

\paragraph{Single language and single source bank.} All items derive from one Vietnamese
exam bank. The four transformation families exploit conventions of Vietnamese exam writing;
their yield in other languages is untested, though the negation and ``all of the above''
patterns are cross-linguistically common.

\paragraph{The transformations assume the source is internally consistent.} The
\textsc{negation} family assumes that if a stem asks for the false option, the other three
are true. Many exam items in practice ask for the \emph{most} false option. This assumption
is what round 2 is designed to test, and it is why \textsc{negation} has the largest audit
rejection count.

\section{Conclusion}

We have shown that multiple-answer medical benchmarks can be built at scale without either
expert authoring or model-generated ground truth, by deriving items from single-answer exam
questions whose keys are already known. The resulting Vietnamese benchmark discriminates
across a very wide capability range---1.8\% to 61.5\% exact match---and is far from
saturated, with 20.3\% of items solved by no model we evaluated.

The benchmark surfaces a capability that single-answer evaluation structurally cannot
measure. A model may hold the knowledge required for an item, as GPT-4o-mini demonstrably
does for 90.5\% of the single-answer items, and still fail it because it cannot decide how
many options to commit to. That deficit is not repaired by giving the model more time to
think, and it improves markedly with model generation, which suggests it should be treated
as a trainable capability and measured accordingly.

\subsection*{Reproducibility statement}

Section~\ref{sec:construction} describes every construction step, including the
deduplication criteria, the mechanical filters applied before any model call, the two-round
LLM protocol and the two audit failure modes we corrected for. Section~\ref{sec:setup}
states the decoding configuration, the output-constraint schema, and the reasoning
condition under which each model was scored, together with the models we could not evaluate
under matched conditions and why. The prompts for both LLM rounds and their JSON schemas
are reproduced in the appendix. We release both benchmarks, their index-aligned
single-answer controls, the per-item provenance records linking every derived item to its
source, the evaluation harness, and the per-item predictions of all nine models, so that
every table in this paper can be recomputed from released artefacts.

\subsection*{Ethics statement}

The source exam bank is licensed material that is not ours to redistribute; what we release
is the derived benchmark, the harness, and per-item model predictions, together with
provenance records sufficient to audit any individual item. The benchmark contains medical
examination questions and no patient data, and no human subjects were involved.

We note two risks. First, the benchmark inherits its ground truth from the source bank and
therefore inherits its errors; we document a concrete instance in the Limitations section
and flag the residual rate rather than claiming a clean key. Scores on this benchmark
should not be read as evidence of clinical safety, and the benchmark is not designed or
validated for deployment decisions about clinical tools. Second, to keep the benchmark
uncontaminated we avoided provider free tiers whose terms permit training on submitted
prompts.

\subsection*{AI use statement}

Generative AI tools were used in two distinct roles in this work, which we separate
deliberately because the distinction is methodologically load-bearing.

Within the benchmark construction pipeline, a language model (Qwen3-32B-FP8) screens
candidate items, rewrites negated stems into affirmative form, writes distractor text, and
performs the round-2 audit, as described in Section~\ref{sec:construction}. The model is
structurally prevented from selecting correct answers: every answer key is derived by rule
from the source item's key, and the model's permitted edits are restricted to removal and
to distractor authorship. This constraint is the central methodological claim of the paper
and is enforced in code rather than by prompting.

Generative AI tools were also used to assist with software implementation and with drafting
and editing the text of this paper. All numerical results reported here were computed by
scripts operating directly on the released prediction files, and each table was verified
programmatically against those files rather than transcribed by hand. We have reviewed all
AI-assisted work and take responsibility for the final content of this paper, including its
text, claims and released artefacts.

\bibliography{iclr2027_conference}
\bibliographystyle{iclr2027_conference}

\appendix
\section{Appendix}

\subsection{Per-family accuracy on Set B}

\begin{table}[h]
\caption{Exact set match (\%) on Set B by transformation family.}
\label{tab:family-acc}
\begin{center}
\begin{small}
\begin{tabular}{lrrrrr}
\toprule
\textbf{Model} & \textsc{negation} & \textsc{meta-all} & \textsc{pool} &
\textsc{meta-letters} & \textbf{single} \\
& (4,573) & (1,092) & (2,810) & (893) & (2,000) \\
\midrule
Gemma3-1B          &  0.2 &  0.0 &  5.1 &  6.0 &  0.1 \\
Qwen3-0.6B         &  1.3 &  2.3 &  5.8 & 11.9 &  1.2 \\
Gemma3-4B          & 20.8 & 28.4 &  9.5 & 25.5 &  4.7 \\
Qwen3-4B           & 20.2 & 38.7 & 16.9 & 38.2 & 12.5 \\
GPT-4o-mini        & 41.0 & 60.4 & 22.0 & 48.2 & 13.2 \\
DeepSeek-v4-Flash  & 43.5 & 60.7 & 33.4 & 56.3 & 41.4 \\
Gemini-2.5-Flash   & 49.7 & 67.2 & 39.9 & 60.1 & 50.9 \\
Claude-Sonnet-5    & 56.6 & 74.0 & 48.0 & 64.6 & 63.6 \\
Gemini-3.6-Flash   & 58.9 & 72.7 & 51.4 & 66.6 & 73.2 \\
\bottomrule
\end{tabular}
\end{small}
\end{center}
\end{table}

The \textsc{pool} family is consistently the hardest, and the single-answer items are the
family on which models differ most: the spread between Gemini-2.5-Flash and
Gemini-3.6-Flash is 22.3 points there against 9.2 on \textsc{negation}. This is the
cardinality effect of Section~\ref{sec:cardinality} appearing per-family.

\subsection{Prompts and schemas}

The full text of the round-1 family-specific prompts, the round-2 audit prompt, and their
JSON schemas are released with the code.

\end{document}
